\section{Experiment 3: Mystery Tubes}

Flooding disrupted the first workday, so this experiment followed a modified, qualitative format. We compared the visible patterns of tubes B and C with the candidate gases. We did not collect quantitative spectra for this part, so the identifications below use the recorded color patterns.

\begin{table}[H]
\centering
\begin{tabular}{>{\raggedright\arraybackslash}p{.08\linewidth}>{\raggedright\arraybackslash}p{.43\linewidth}>{\raggedright\arraybackslash}p{.36\linewidth}}
\toprule
Tube & Recorded features & Best qualitative match \\
\midrule
B & Prominent red, orange, and yellow features. & Neon, whose visible spectrum is rich in orange and red lines. \\
C & Many lines, including blue, two distinct green regions, and yellow to red. & Argon, whose lines span a broad range of visible colors. \\
\bottomrule
\end{tabular}
\caption{Observations and candidate identifications. Line positions and relative intensities were not measured.}
\end{table}

Tube B's concentration of prominent warm colored features is consistent with neon. Tube C showed a more varied pattern across the visible range, which is consistent with argon. These comparisons use the handout's descriptions and the \href{https://physics.nist.gov/PhysRefData/Handbook/Tables/neontable2.htm}{NIST neon} and \href{https://physics.nist.gov/PhysRefData/Handbook/Tables/argontable2.htm}{argon line tables}.

The third candidate, molecular nitrogen, can also produce a complex spectrum. Its rotational and vibrational structure creates groups of closely spaced transitions that appear as bands. The notes do not resolve that band structure, so the argon identification is less secure than a match using measured wavelengths.

The spectral pattern carries more information than the overall glow of a tube. A discharge excites the gas, and emitted photon energies reflect differences between its energy levels,
\begin{equation}
\frac{hc}{\lambda}=E_i-E_f.
\end{equation}
Different gases have different energy spacings, so the combination of spectral features acts as a fingerprint. Our observations favor B as neon and C as argon, with the confidence limited by the qualitative comparison.
